\scanpage{37}
% L1
\uuline{longueur finie}. Elle est stable par $\otimes$, sur cette catégorie les $\mathrm{Ext}^i$ $(i>0)$ sont nuls, et la catégorie est aussi stable par $\mathscr{H}\mathrm{om}$. (cas de $M(X)_{\mathrm{spec}}$). On a une ``catégorie tensorielle \unsure{''} sur $\mathbb Q$''.

% L6
Le \deleted{\unsure{la}} $\lambda$-anneau défini par cette catégorie est noté $\unsure{\mathbf M|\mathbb M|M}^+(X)$ resp. $M(X)$. Noter que $\underline{\xi}=\mathrm{classe}\,\deleted{\mathbb Z}\mathbb Q_X(-1))$ est un élément de $M^+(X)$, et on a
\[
M(X)=M^+(X)_\xi=M^+(X)\left[\frac{1}{\xi}\right].
\]

% L11
En fait, $\xi$ est non div. de $0$ : $\mathbf M^+(X)$ comme $\mathbb Z$-module admet comme base les classes d'\deleted{\unsure{éléments}}\inserted{\uline{objets}} \uline{simples} de $M^+(X)$, et $x\mapsto\xi\cdot x$ dans $\mathbf M^+(X)$ est défini par une application \uline{injective} sur l'ens de ses objets simples. [N.B. $M$ simple \inserted{ds $\mathbf M^+(X)$} $\Longleftrightarrow$ \inserted{$M$ simple dans $M(X)$ $\Longleftrightarrow$} $\mathbf M(-1)$ simple, comme il résulte du fait que $\mathbf M^+(X)$ est épaisse dans $M(X)$ et $\otimes\mathbb Z(-1)$ une auto équivalence ]
% archive sheet number: 36
